% !TEX root = ../main.tex

\section{Piano Keys \& Accidentals}

\begin{itemize}
    \item \textbf{Sharp ($\sharp$):} Move 1 half-step (semitone) to the right (higher).
    \item \textbf{Flat ($\flat$):} Move 1 half-step (semitone) to the left (lower).
    \item \textbf{Enharmonics:} The same physical key with two names:
    \begin{itemize}
        \item C$\sharp$ / D$\flat$ (between C and D)
        \item D$\sharp$ / E$\flat$ (between D and E)
        \item F$\sharp$ / G$\flat$ (between F and G)
        \item G$\sharp$ / A$\flat$ (between G and A)
        \item A$\sharp$ / B$\flat$ (between A and B)
    \end{itemize}
\end{itemize}

\section{Scale Formulas}

From the \href{https://www.youtube.com/watch?v=GugYaekR8KQ}{Video: Scale for Beginners}.

\begin{itemize}
    \item \textbf{Major Scale:} Root $\rightarrow$ W $\mid$ W $\mid$ H $\mid$ W $\mid$ W $\mid$ W $\mid$ H
    \item \textbf{Minor Scale:} Root $\rightarrow$ W $\mid$ H $\mid$ W $\mid$ W $\mid$ H $\mid$ W $\mid$ W
    \item $W$ = Whole step (2 semitones / keys).
    \item $H$ = Half step (1 semitone / adjacent key).
\end{itemize}

\begin{warningbox}{The Golden Rule of Scale Spelling}
Standard diatonic scales must use \textbf{every letter (A, B, C, D, E, F, G) exactly once} before repeating the root at the octave. Never skip or duplicate a letter (e.g., $D\flat$ Major has $F$ and $G\flat$, not $F$ and $F\sharp$).
\end{warningbox}

\section{Fingering \& Hand Mechanics}

Using the right finger standardizes muscle memory for speed, accuracy, and reliability.
\begin{itemize}
    \item Fingers: \textbf{1} = Thumb, \textbf{2} = Index, \textbf{3} = Middle, \textbf{4} = Ring, \textbf{5} = Pinky.
    \item \textbf{Thumb on white keys:} Thumbs tuck under to white keys; avoid placing the thumb on black keys in basic scale patterns.
    \item \textbf{Crossing rule:} Every 7-note scale is split into a 3-note group and a 4-note group ($3 + 4 = 7$):
    \begin{itemize}
        \item Entering 3-note group $\rightarrow$ cross finger \textbf{3} over thumb.
        \item Entering 4-note group $\rightarrow$ cross finger \textbf{4} over thumb.
    \end{itemize}
\end{itemize}

\section{Scales Library}

\subsection{D-flat Major}
\begin{definitionbox}{D$\flat$ Major Scale}
\begin{itemize}
    \item \textbf{Notes:} D$\flat$ -- E$\flat$ -- F -- G$\flat$ -- A$\flat$ -- B$\flat$ -- C -- D$\flat$
    \item \textbf{RH Ascending:} 2 -- 3 -- 1 -- 2 -- 3 -- 4 -- 1 -- 2
    \item \textbf{RH Descending:} 2 -- 1 -- 4 -- 3 -- 2 -- 1 -- 3 -- 2
    \item \textbf{LH Ascending:} 3 -- 2 -- 1 -- 4 -- 3 -- 2 -- 1 -- 3
    \item \textbf{LH Descending:} 3 -- 1 -- 2 -- 3 -- 4 -- 1 -- 2 -- 3
\end{itemize}
\end{definitionbox}

\subsection{C Major}
\begin{definitionbox}{C Major Scale (White Keys Baseline)}
\begin{itemize}
    \item \textbf{Notes:} C -- D -- E -- F -- G -- A -- B -- C
    \item \textbf{RH Ascending:} 1 -- 2 -- 3 -- 1 -- 2 -- 3 -- 4 -- 5
    \item \textbf{RH Descending:} 5 -- 4 -- 3 -- 2 -- 1 -- 3 -- 2 -- 1
    \item \textbf{LH Ascending:} 5 -- 4 -- 3 -- 2 -- 1 -- 3 -- 2 -- 1
    \item \textbf{LH Descending:} 1 -- 2 -- 3 -- 1 -- 2 -- 3 -- 4 -- 5
\end{itemize}
\end{definitionbox}

\subsection{Future Scales (Queue)}
\begin{itemize}
    \item \textbf{Upcoming:} G Major, F Major, A minor.
\end{itemize}
% Template for adding new scales:
% \begin{definitionbox}{[Scale Name]}
% \begin{itemize}
%     \item \textbf{Notes:} ...
%     \item \textbf{RH Ascending:} ...
%     \item \textbf{RH Descending:} ...
%     \item \textbf{LH Ascending:} ...
%     \item \textbf{LH Descending:} ...
% \end{itemize}
% \end{definitionbox}
